\section{The actual standard and dual Bessel quotients}
\label{sec:bessel-quotient-suppliers-v52}

This chapter supplies the two original \texttt{QuotPStatement} and
\texttt{QuotMStatement}.  The spaces at the end are images of the actual
tensor coefficient maps, and the kernel being removed is the kernel of
all of the prescribed periods.  The intermediate statements retain the
broader source scopes: arbitrary generated-module tuples, arbitrary
cyclic lowest-weight modules, and intertwining maps into inner-product
spaces which need not be complete.

We use the reviewed DC, BSP, BST, BSE and BNV chapters.  Coordinates are
$0,\ldots,n-1$, $G_\epsilon=U(J_\epsilon)$ and
$K_\epsilon=G_\epsilon\cap U(n)$.  Write $p=p(\epsilon)$, $q=q(\epsilon)$,
$q_0=\min(p,q)$,
\[
 \rho_i=(n-1)/2-i,\qquad
 \Lambda(\lambda)=\lambda+\rho_n-\rho_c,\qquad
 c(\lambda)=\|\lambda\|^2-\|\rho\|^2.
 \tag{BQS.1}
\]
Positive roots are $\alpha_{ab}=e_a-e_b$ for $a<b$; compact roots have
$\epsilon_a=\epsilon_b$.  The \emph{integral positive-root cone} is
\[
 Q_+=\left\{\sum_{a<b}m_{ab}\alpha_{ab}:m_{ab}\in\mathbb Z_{\ge0}\right\}.
 \tag{BQS.2}
\]
For a weight $\gamma$ put $P_j(\gamma)=\sum_{i<j}\gamma_i$ and
$\operatorname{ht}(\gamma)=\sum_i(n-1-i)\gamma_i$.
Every $\gamma\in Q_+$ has integral coordinates, nonnegative integral
partial sums, total sum zero, and
$\operatorname{ht}(\gamma)=\sum_{a<b}(b-a)m_{ab}$.
Consequently $Q_+\cap(-Q_+)=\{0\}$, and a nonempty sum of positive
roots has strictly positive height.

An \emph{original unitary datum} below means every datum quantified in
\texttt{LowestKTypeEuler\allowbreak{}TensorStandard}: a complete complex Hilbert space
$H$, an actual continuous unitary representation $\sigma$ of $G_\epsilon$,
a derived invariant domain $D\subset H$ with matrix-unit operators
$E_{ab}$, and a finite-dimensional compact irreducible
$V\subset D$ of highest weight $\Lambda(\lambda)$, having Casimir
$c(\lambda)$ and stable under the actual compact group.  The parameter is
strictly decreasing and
$\lambda_i-(n-1)/2\in\mathbb Z$.  Also $p,q\ge1$ and
$\operatorname{GapLap}(\epsilon,\lambda,h)$ holds: the noncompact gaps
are at least $h$ and $h>h_{\rm need}(\epsilon)$, where exactly
\[
\begin{split}
h_{\rm need}(\epsilon)=\max\bigl\{&
2+\max(3n/2,1+npq),\
16nq_0+8\sqrt{2q_0}+2,\\
&8\sqrt{2q_0}\,[4\max(2nq_0,2(q-1)n^2)+1]\bigr\}.
\end{split}
\tag{BQS.3}
\]
In particular this datum is not assumed to come with a chosen coefficient
identity.  BSP2 supplies that identity at this exact gap when it is used.
Put $S=U(\mathfrak{gl}_n)V$.  The notation \emph{compact highest vector
of weight $\nu$} means $E_{aa}v=\nu_av$ and $E_{ab}v=0$ for
$a<b$ of equal sign; a zero vector is allowed in this predicate.

For any complex $\mathfrak{gl}_n$-module $M$, use the literal tensor
operators
\[
\begin{aligned}
 (A^+_{ab}f)_d&=E_{ab}f_d+\delta_{da}f_b,&
 (\Omega^+f)_d&=\sum_cE_{cd}f_c,\\
 (A^-_{ab}f)_d&=E_{ab}f_d-\delta_{db}f_a,&
 (\Omega^-f)_d&=-\sum_cE_{dc}f_c,\\
 \theta_s^+&=\lambda_s-(n-1)/2,&
 \theta_s^-&=-\lambda_s-(n-1)/2,\\
 X_s^\pm&=S^n\cap\ker(\Omega^\pm-\theta_s^\pm).
\end{aligned}
\tag{BQS.4}
\]
These are \texttt{tensV}, \texttt{tensVd}, \texttt{eeV},
\texttt{eeVd} and \texttt{Xs}; products of operators act from right to
left.

% BEGIN BQS1
\paragraph{BQS1. Tensor generation, commuting Euler operators and finite compact hulls.}
\label{claim:BQS1}
For every complex module $M$, every matrix-unit representation $E$ and
every subspace $V$, the module $S^n$, $S=U(\mathfrak{gl}_n)V$, is generated
by $V^n$ under either $A^+$ or $A^-$.  Both tensor actions satisfy the
matrix-unit relations and commute with the corresponding $\Omega^\pm$.
If $V$ is finite-dimensional and stable under the compact matrix units,
every $f\in X_s^\pm$ is contained in a finite-dimensional compact-stable
subspace of $X_s^\pm$.  No Hilbert completion is required.

\begin{proof}
Write $[v]_c$ for the tuple with $v$ in coordinate $c$ and zero elsewhere.
The identities
\[
 A^+_{ab}[v]_c=[E_{ab}v]_c+\delta_{bc}[v]_a,\qquad
 A^-_{ab}[v]_c=[E_{ab}v]_c-\delta_{ac}[v]_b
 \tag{BQS.5}
\]
show inductively that the tensor module generated by $V^n$ contains
$[E_{a_1b_1}\cdots E_{a_mb_m}v]_c$ for every word and $v\in V$.
The induction subtracts the second term of (BQS.5), whose word length
is smaller.  The reverse inclusion follows since both terms preserve
$S^n$.  This proves the full original \texttt{GenStatement}.

The tensor relations follow either by substituting (BQS.5) or by adding
the commuting actions on the two factors $M\otimes\mathbb C^n$ and
$M\otimes(\mathbb C^n)^*$.  To check Euler commutation in the standard
case, the $d$-component of $[\Omega^+,A^+_{ab}]f$ is
\[
 \sum_c[E_{cd},E_{ab}]f_c+E_{ad}f_b
       -\delta_{da}\sum_cE_{cb}f_c=0.
 \tag{BQS.6}
\]
Indeed $[E_{cd},E_{ab}]=\delta_{da}E_{cb}-\delta_{bc}E_{ad}$.
For the dual case the same substitution gives
$-\sum_c[E_{dc},E_{ab}]f_c+E_{db}f_a
-\delta_{db}\sum_cE_{ac}f_c=0$.
Thus each Euler eigenspace is a full tensor submodule.

Let $F_m$ be the span of all words of length at most $m$ applied to $V$.
It is finite-dimensional.  A compact matrix unit preserves it: move that
unit past each word letter using the commutator relation, and then use
compact stability of $V$.  Every tuple in $S^n$ lies in some $F_m^n$.
The intersection $F_m^n\cap X_s^\pm$ is compact-stable, because
$F_m^n$ is stable under the compact tensor action and $\Omega^\pm$
commutes with that action.  It is the required finite hull.
\end{proof}
% END BQS1

% BEGIN BQS2
\paragraph{BQS2. The adjugate identities give the whole Euler polynomial.}
\label{claim:BQS2}
Suppose $S$ is an invariant subspace of a complex matrix-unit module and
the literal ordered Capelli operator has scalar value
\[
 \mathcal C_E(t)|_S=p(t)I,\qquad
 p(t)=\prod_i(t+\theta_i),\qquad t\in\mathbb C.
 \tag{BQS.7}
\]
Then on every tuple in $S^n$,
\[
 \prod_i(\Omega^+-\theta_i)=0,\qquad
 \prod_i\bigl(\Omega^-+\theta_i+n-1\bigr)=0.
 \tag{BQS.8}
\]
This assertion needs neither distinct roots nor an inner product.

\begin{proof}
For $n=0$, $S^0$ is the zero vector space, so its identity is zero.
Assume $n\ge1$.  BST8 gives, with exactly its ordered complementary
matrices $U(t),V(t)$,
\[
 (\Omega^++tI)U(t)=p(t)I,\qquad
 \bigl(-\Omega^-+(t-n+1)I\bigr)V(t)=p(t)I
 \quad\hbox{on }S^n.
 \tag{BQS.9}
\]
Here the matrices of operators act on column tuples.  Each minor
preserves $S$, so these are genuine identities of endomorphism-valued
polynomials on $S^n$.  Their scalar validity for every $t$ implies
coefficient validity: apply an arbitrary linear functional to their
value on an arbitrary vector and use the scalar polynomial identity.
The diagonal complementary minor has leading term $t^{n-1}$; every
off-diagonal complementary minor has degree at most $n-2$.  Thus both
$U$ and $V$ are monic of degree $n-1$ as operator polynomials.

For clarity, the elementary substitution argument in a noncommutative
endomorphism algebra is as follows.  If
$(tI-B)(t^{n-1}I+Q_1t^{n-2}+\cdots+Q_{n-1})=p(t)I$,
comparison from the highest coefficient downward gives
$Q_j=BQ_{j-1}+p_jI$, where
$p(t)=t^n+p_1t^{n-1}+\cdots+p_n$ and $Q_0=I$.
The constant coefficient then gives
$B^n+p_1B^{n-1}+\cdots+p_nI=0$.
This proves $p(B)=0$ without substituting into polynomials with
uncontrolled noncommuting coefficients.
Apply it first to $B=-\Omega^+$ and then to
$B=\Omega^-+(n-1)I$ in (BQS.9).
Removing the harmless scalar $(-1)^n$ in the first case gives (BQS.8).
\end{proof}
% END BQS2

% BEGIN BQS3
\paragraph{BQS3. Both original lowest-type Euler statements and their projectors.}
\label{claim:BQS3}
For every original unitary datum and every $f\in S^n$,
\[
 \prod_i(\Omega^+-\theta_i^+)f=0,\qquad
 \prod_i(\Omega^--\theta_i^-)f=0.
 \tag{BQS.10}
\]
These are the full original \texttt{LowestKTypeEuler\allowbreak{}TensorStandard}
and \texttt{LowestKTypeEuler\allowbreak{}TensorDual}.  For either sign the operators
\[
 \Pi_s^\pm=
 \prod_{t\ne s}\frac{\Omega^\pm-\theta_t^\pm}
                       {\theta_s^\pm-\theta_t^\pm}
 \quad\hbox{on }S^n
 \tag{BQS.11}
\]
are commuting idempotents with pairwise zero products, sum $I$, and
image $X_s^\pm$.  They commute with every corresponding tensor matrix
unit, not only the compact ones.

\begin{proof}
The gap (BQS.3) permits the full cyclic rigidity assertion of BSP2.
Choose its actual compact isometry and coefficient identity for this
datum.  BSP3, applied to that identity, gives
$p(t)=\prod_i(t+\lambda_i-(n-1)/2)$ on all of $S$.
BQS2 yields both assertions in (BQS.10), because
$-\theta_i^+-(n-1)=\theta_i^-$.
Thus no coefficient identity has been added to the hypotheses of
either original Euler statement.

The numbers $\theta_i^\pm$ are distinct by strict decrease of $\lambda$.
The scalar Lagrange polynomials $l_s$ satisfy
$l_s(\theta_t)=\delta_{st}$; the polynomials
$\sum_s l_s-1$, $l_s^2-l_s$, $l_sl_t$ for $s\ne t$, and
$(z-\theta_s)l_s(z)$ are divisible by
$\prod_i(z-\theta_i)$.  Evaluate these scalar polynomial identities
at $\Omega^\pm$ and use (BQS.10).
They give the stated projections and sum.  An eigenvector of eigenvalue
$\theta_s$ is fixed by $\Pi_s$, proving equality of its image with
$X_s^\pm$.  Commutation follows from BQS1.
\end{proof}
% END BQS3

% BEGIN BQS4
\paragraph{BQS4. The whole polynomial identity for a cyclic lowest-weight module.}
\label{claim:BQS4}
Let $M$ be any complex matrix-unit module generated by $\phi$, and suppose
\[
 E_{ab}\phi=0\quad(b<a),\qquad
 E_{aa}\phi=(\lambda_a+\rho_a)\phi.
 \tag{BQS.12}
\]
No regularity, integrality, positivity or strict decrease of $\lambda$
is assumed.  Both products in (BQS.10) vanish on every tuple of $M^n$.
This is the full original \texttt{VermaProductIdentity}.

\begin{proof}
Set $v_t=[\phi]_t$.  Directly from (BQS.4), (BQS.5) and (BQS.12),
\[
\begin{aligned}
 (\Omega^+-\theta_t^+)v_t
    &=\sum_{d>t}A^+_{td}v_d,\\
 (\Omega^--\theta_t^-)v_t
    &=-\sum_{d<t}A^-_{dt}v_d .
\end{aligned}
\tag{BQS.13}
\]
For example the $t$-coordinate in the first identity is
$(\lambda_t+\rho_t-\theta_t^+)\phi=(n-1-t)\phi$;
the right side has exactly $n-1-t$ such contributions.  The other
coordinates are $E_{td}\phi$ for $d>t$, and zero for $d<t$.
In the second identity the $t$-coordinate is $t\phi$, and the same
coordinate check gives the displayed minus sign.

For the standard action let $F_j$ be the tensor submodule generated by
$v_t$, $t\ge j$.  For the dual action let $F_j$ be generated by
$v_{n-1-t}$, $t\ge j$.  In either case BQS1 gives $F_0=M^n$ and
$F_n=0$.  The $j$th quotient is cyclic, generated by the corresponding
$v_t$, and (BQS.13) says that $\Omega$ has scalar value $\theta_t$
on it.  BQS1 allows this equality to propagate to the whole quotient.
Hence $(\Omega-\theta_t)F_j\subset F_{j+1}$.
The product of the $n$ factors sends $F_0$ into $F_n=0$.
All factors are polynomials in one operator, so their order is
immaterial.  This proves both signs, including repeated roots.
\end{proof}
% END BQS4

% BEGIN BQS5
\paragraph{BQS5. The full cyclic-module cone and extremal-line theorem.}
\label{claim:BQS5}
Under (BQS.12), now assume $\lambda$ strictly decreasing.  Every
nonzero simultaneous diagonal weight vector $f$ of weight $\nu$ in
$\ker(\Omega^\pm-\theta_s^\pm)\subset M^n$ satisfies
\[
              \nu-(\lambda+\rho\pm e_s)\in Q_+.
 \tag{BQS.14}
\]
The simultaneous weight space of weight $\lambda+\rho\pm e_s$ in
that eigenspace has dimension at most one.  This is the whole original
\texttt{VermaConeStatement}; it does not require $f$ to be highest
for any compact group.

\begin{proof}
Use the flag from BQS4, writing its ordered eigenvalues as
$\vartheta_0,\ldots,\vartheta_{n-1}$ and its $j$th generator as $v$.
Here $j=s$ in the standard case and $j=n-1-s$ in the dual case.
All $\vartheta_i$ are distinct.  If $\Omega f=\vartheta_j f$,
applying the first $j$ flag factors shows
$\prod_{i<j}(\vartheta_j-\vartheta_i)f\in F_j$.
Its nonzero scalar coefficient implies $f\in F_j$.

Modulo $F_{j+1}$ every lower matrix unit kills $v$:
in the standard case $A^+_{ab}v_s=\delta_{bs}v_a$ for $a>b$,
and $a>s$ puts this in the next flag term; in the dual case
$A^-_{ab}v_s=-\delta_{as}v_b$ and $b<s$ does the same.
The diagonal weight of $v$ is $b=\lambda+\rho\pm e_s$.
Commuting lower units to the right and diagonals past upper units
therefore shows that $F_j/F_{j+1}$ is spanned by positive-root
words applied to $v$.  This is only PBW spanning; no freeness or
irreducibility of the cyclic module is used.

Put $R=\prod_{i>j}(\Omega-\vartheta_i)$ and
$a=\prod_{i>j}(\vartheta_j-\vartheta_i)\ne0$.
The flag argument gives $R F_{j+1}=0$.  Thus every such eigenvector
$f$ satisfies
\[
           a f=R y
 \quad\hbox{for a finite sum $y$ of positive-root words on $v$}.
 \tag{BQS.15}
\]
The operator $R$ commutes with all tensor matrix units.  Each word of
length $m$ on $v$ has diagonal weight $b+\beta$, where
$\beta\in Q_+$ and $\operatorname{ht}(\beta)\ge m$.
Its image under $R$ has the same weight or is zero.

Distinct simultaneous weight spaces are linearly independent, even in
an arbitrary module: for any finite set of distinct weights choose a
linear combination of the diagonal operators separating them, and use
its scalar Lagrange polynomials.  Apply this observation to the finite
sum (BQS.15).  A nonzero weight vector can therefore have only a weight
$b+\beta$, proving (BQS.14).
For weight exactly $b$, every positive-length word has a different
weight by the strict positivity of height.  Its component in (BQS.15)
is consequently a scalar multiple of $Rv$.
The extremal weight space is contained in this one line, which may
itself be zero.  This proves the asserted dimension bound for both signs.
\end{proof}
% END BQS5

% BEGIN BQS6
\paragraph{BQS6. The actual principal-minor cyclic module.}
\label{claim:BQS6}
Let $\mu=\lambda+\rho\in\mathbb Z^n$, and let $\Phi_\mu$ be the
principal-minor function of DC2.  Write
\[
 \Omega_{\rm hol}=\{z:\Delta_j(z)\ne0\ (1\le j\le n)\},\quad
 \Omega_P=\{z:z^{-1}\in\Omega_{\rm hol}\},\quad
 \Phi_P(z)=\Phi_\mu(z^{-1}).
 \tag{BQS.16}
\]
On holomorphic functions on $\Omega_P$, extended by zero outside it,
define
$\mathcal E_{ab}F(z)=dF_z(zE_{ab})$ on $\Omega_P$ and zero outside.
The cyclic module $M_\mu=U(\mathfrak{gl}_n)\Phi_P$ is a literal module
of such functions.  Its generator satisfies (BQS.12), and it is
generated by that vector.  These are all fields of the original
\texttt{k2vM} model.

\begin{proof}
All leading minors are polynomial, and the final one is the determinant.
Thus $\Omega_{\rm hol}$ and $\Omega_P$ are open subsets of invertible
matrices.  Integer powers of their nonzero minors are holomorphic.
For a holomorphic $F$, $z\mapsto dF_z(zX)$ is holomorphic there;
extension by zero is only a function outside this open set and is not
claimed to be smooth across its boundary.
Commuting the two vector fields $z\mapsto zX,z\mapsto zY$ gives
the field $z\mapsto z[X,Y]$, so the $\mathcal E_{ab}$ obey the
matrix-unit relations, also outside $\Omega_P$ where they all vanish.

For $a>b$, $(z\exp(tE_{ab}))^{-1}=\exp(-tE_{ab})z^{-1}$.
Triangular multiplication in DC2 leaves every leading minor unchanged
under this lower unipotent multiplication.  Its derivative is zero.
For $E_{aa}$ the same identity and DC2's diagonal character give
$\Phi_P(z\exp(tE_{aa}))=e^{t\mu_a}\Phi_P(z)$ locally.
Thus $\mathcal E_{aa}\Phi_P=\mu_a\Phi_P$ and lower units kill the
generator.  Restricting the operators to their generated submodule
preserves all identities and makes generation tautological.
This is exactly the source's zero-extended \texttt{MMT}, with no
unproved global analytic continuation attached to it.
\end{proof}
% END BQS6

\subsection{The compact integral transfers the cone to the actual datum}

% BEGIN BQS7
\paragraph{BQS7. Full transfer of finite derivative relations through the compact average.}
\label{claim:BQS7}
Let $\lambda$ be Harish--Chandra and let $T$ be any original
compact type of weight $\Lambda(\lambda)$,
$w_T$ its unit highest vector, and $\mu=\lambda+\rho$.
For arbitrary finitely many lists $\mathbf Z_i$ of complex matrices
and arbitrary $v_i\in T$, suppose
\[
 \sum_i\langle T(k)w_T,v_i\rangle
       R^{\rm hol}_{\operatorname{Ad}(k^{-1})\mathbf Z_i}
                         \Phi_P(z)=0
 \quad(k\in K_\epsilon,\ z\in\Omega_P).
 \tag{BQS.17}
\]
Then, for every $g\in G_\epsilon$,
\[
                  \sum_i R_{\mathbf Z_i}\Psi_{T,\mu}(g)v_i=0.
 \tag{BQS.18}
\]
Here every list keeps its order, $R^{\rm hol}_ZF(z)=dF_z(zZ)$, and
$R_Z$ is the original complexified tangential right derivative.
Empty lists and zero summands are allowed.  This is the full original
\texttt{TransferStatement}.

\begin{proof}
Use probability Haar on $K$; replacing it by the source's Haar changes
the coefficient $\dim T$ into $\dim T/\mu_K(K)$ and changes no identity.
DC2 gives a common open neighborhood $\mathcal U$ of the positive cone
on which every $k^{-1}zk$ is in $\Omega_{\rm hol}$ and
\[
 P(z)=(\dim T)\int_K\Phi_\mu(k^{-1}zk)
                 T(k)P_{w_T}T(k)^{-1}\,dk.
 \tag{BQS.19}
\]
All derivatives of this identity are obtained under the integral.
Indeed for a fixed $z\in\mathcal U$ choose a smaller relatively compact
complex neighborhood whose closure remains in $\mathcal U$.
The conjugated minors have a uniform positive lower bound in modulus
on its product with compact $K$.  Every fixed derivative is a continuous
rational expression with these nonzero denominators, so it is uniformly
bounded there.  Differentiation under this finite measure integral follows.
This argument works for every finite word, without one word-length bound
being imposed on the statement.

Fix $g$.  DC3 supplies $k_0\in K$ and a neighborhood on which
$\Psi(g')=T(k_0)P((g')^{-1}k_0)$.  Put $z_0=g^{-1}k_0$.
If $L_ZF(z)=dF_z(Zz)$, the inverse rule gives
\[
 R_{\mathbf Z}\Psi(g)
       =T(k_0)L_{-\mathbf Z}P(z_0),
 \qquad
 R^{\rm hol}_{\mathbf Z}\Phi_P(y)
       =L_{-\mathbf Z}\Phi_\mu(y^{-1}).
 \tag{BQS.20}
\]
Induction proves these formulas with the displayed order: at each
step $(g'\exp(tZ))^{-1}=\exp(-tZ)(g')^{-1}$.
Complex linearity gives the same result for every complex matrix $Z$.

Differentiate (BQS.19), apply it to each $v_i$, and sum.
The rank-one factor satisfies
\[
 T(k)P_{w_T}T(k)^{-1}v_i
             =\langle T(k)w_T,v_i\rangle T(k)w_T.
 \tag{BQS.21}
\]
At this $k$ use (BQS.20) with
$y=(k^{-1}z_0k)^{-1}\in\Omega_P$.
The scalar multiplying $T(k)w_T$ is exactly the sum in (BQS.17).
It is zero at every $k$.  The integral and hence its left multiplication
by $T(k_0)$ are zero.  This proves (BQS.18).
\end{proof}
% END BQS7

% BEGIN BQS8
\paragraph{BQS8. Formal word space, compact action and highest-vector lifting.}
\label{claim:BQS8}
Let $D$ be any complex inner-product module, possibly incomplete, on
which $E$ satisfies the matrix-unit relations and
$E_{ab}^*=E_{ba}$ for compact pairs.  Let $V$ be finite-dimensional
and compact-stable.  Form the algebraic direct sum
\[
 \mathcal W=\bigoplus_{m\ge0}
             \mathfrak{gl}_n^{\otimes m}\otimes V\otimes\mathbb C^n .
 \tag{BQS.22}
\]
Its elementary vectors are denoted $[\mathbf Z,v,c]$; expansion in
matrix units identifies it with the original finitely supported word
space \texttt{KS}.  Define
\[
\begin{aligned}
 \mathcal T[\mathbf Z,v,c]&=[E_{\mathbf Z}v]_c,\\
 \mathcal B[\mathbf Z,v,c]&=\sum_d[(E_{cd},\mathbf Z),v,d],\\
 \mathcal A_X[\mathbf Z,v,c]&=
 [\operatorname{ad}(X)\mathbf Z,v,c]+[\mathbf Z,E_Xv,c]
                 +[\mathbf Z,v,Xe_c]\quad(X\in\mathfrak k_{\mathbb C}).
\end{aligned}
\tag{BQS.23}
\]
The adjoint term means the sum over the list's individual positions.
Then $\mathcal T$ maps onto $S^n$,
$\mathcal T\mathcal B=\Omega^+\mathcal T$,
$\mathcal T\mathcal A_{ab}=A^+_{ab}\mathcal T$ for compact pairs, and
$[\mathcal A_{ab},\mathcal B]=0$.
Every compact highest vector $f\in S^n$ of weight $\nu$ has a compact
highest lift $u\in\mathcal W$ of that same weight, with $\mathcal T u=f$.
This includes the full original \texttt{lift\_highestWeightVector}.

\begin{proof}
Surjectivity follows by expanding each component of $f$ as a finite sum
of words on $V$.  Prepending the indicated matrix unit proves
$\mathcal T\mathcal B=\Omega^+\mathcal T$.
For compact $X$, the Leibniz commutator identity
$[E_X,E_{\mathbf Z}]=E_{\operatorname{ad}(X)\mathbf Z}$
proves the second intertwining identity.
The $\mathcal A$ are the sum of the adjoint, $V$, and standard tensor
actions, hence satisfy all compact matrix-unit relations.
The element $\sum_{c,d}E_{cd}\otimes E_{dc}$ is invariant under
simultaneous conjugation.  Expanding its commutator with a matrix unit
cancels the two Kronecker sums; this is exactly
$[\mathcal A_{ab},\mathcal B]=0$ on (BQS.23).

Give each matrix-letter space its Hilbert--Schmidt inner product, give
$V$ its inherited inner product, and take the tensor/direct-sum inner
products in (BQS.22).  On compact pairs,
$(\operatorname{ad}E_{ab})^*=\operatorname{ad}E_{ba}$,
as follows from $\operatorname{tr}(X^*Y)$ and cyclicity of trace.
The same adjoint identity holds on the other two factors, and therefore
$\mathcal A_{ab}^*=\mathcal A_{ba}$.
Words of length at most $m$ form a finite-dimensional compact-stable
space $\mathcal W_{\le m}$.

Choose any preimage of $f$ in some $\mathcal W_{\le m}$.
The kernel of $\mathcal T$ on that space is compact-stable.  Its
orthogonal complement is also compact-stable by the adjoint identities.
The restriction of $\mathcal T$ to this complement is an isomorphism
onto $\mathcal T(\mathcal W_{\le m})$.
Lift $f$ by its inverse.  The intertwining identities and injectivity
show that the lift has the same diagonal eigenvalues and is killed
by every compact upper unit.  All orthogonal complements used here
are finite-dimensional; completeness of $D$ was never invoked.
\end{proof}
% END BQS8

% BEGIN BQS9
\paragraph{BQS9. The full trace-kernel implication for the actual coefficient.}
\label{claim:BQS9}
Let $\lambda$ be Harish--Chandra, and let $(H,\sigma,D,E,V)$
have the complete unitary, derived and compact irreducibility
properties of an original datum.  Suppose an original
$T,\mu$ and linear isomorphism $\iota:V\to T$ satisfy the exact coefficient
identity
\[
 \langle v',\sigma(g)v\rangle
       =\langle\iota v',\Psi_{T,\mu}(g)\iota v\rangle
       \quad(g\in G_\epsilon,\ v,v'\in V).
 \tag{BQS.24}
\]
For this assertion no gap or Casimir hypothesis is additionally required.
Put $v_0=\iota^{-1}w_T$, $\ell(v)=\langle v_0,v\rangle$, and
\[
 \Theta[\mathbf Z,v,c]
       =[\ell(v)\mathcal E_{\mathbf Z}\Phi_P]_c
                   \quad\hbox{in }M_\mu^n.
 \tag{BQS.25}
\]
For \emph{every} compact-stable subspace $W\subset\mathcal W$,
\[
                    \Theta W=0\quad\Longrightarrow\quad\mathcal T W=0.
 \tag{BQS.26}
\]
This proves the full original \texttt{lowestKType\_trace\_integral},
and hence \texttt{lowestKType\_trace}.

\begin{proof}
At $g=1$, (BQS.24) and $\Psi(1)=1$ show that $\iota$ is an isometry.
The finite-dimensional compact-stable space $V$ is stable under the
actual compact group.  To see this directly, solve the orbit differential
equation for each compact real Lie direction on $V$; its matrix
exponential agrees with the actual $\sigma$-orbit by uniqueness.
Every unitary block matrix is the exponential of a skew-Hermitian block
matrix, so this covers $K$.
Putting $g=k$ in (BQS.24), using $\Psi(k)=T(k)$ and all $v'\in V$,
now gives $\iota\sigma(k)|_V=T(k)\iota$.

The actual compact action on the word space is
\[
 \mathcal U(k)=
  \bigoplus_m\operatorname{Ad}(k)^{\otimes m}
                    \otimes\sigma(k)|_V\otimes k.
 \tag{BQS.27}
\]
Its derivative is $\mathcal A$.  It preserves word length.
For each $x\in W$, choose $m$ with $x\in\mathcal W_{\le m}$.
The intersection $W\cap\mathcal W_{\le m}$ is finite-dimensional and
$\mathcal A$-stable.  Matrix exponentials and the preceding logarithms
therefore show $\mathcal U(k)x\in W$ for every $k$.
Consequently $\Theta(\mathcal U(k^{-1})x)=0$.

Write $x=\sum_i[\mathbf Z_i,v_i,c_i]$.
Evaluation of the last zero function at $z\in\Omega_P$ gives, in
coordinate $d$,
\[
 \sum_i(k^{-1})_{dc_i}\,
   \langle T(k)w_T,\iota v_i\rangle
   R^{\rm hol}_{\operatorname{Ad}(k^{-1})\mathbf Z_i}
                                 \Phi_P(z)=0.
 \tag{BQS.28}
\]
Here unitarity moved $\sigma(k^{-1})$ from the second slot to the
first slot of $\ell$; the isometry $\iota$ then gave the displayed
coefficient.  Multiplication of these coordinate equations by $k$
isolates each fixed index $c$ and gives the hypothesis of BQS7 for
the summands with $c_i=c$.
It follows that
$\sum_{c_i=c}R_{\mathbf Z_i}\Psi(g)\iota v_i=0$
for all $g$.

Differentiate (BQS.24) through these finite words.  The derived
domain is invariant, so induction gives
\[
 \left\langle v',\sigma(g)
       \sum_{c_i=c}E_{\mathbf Z_i}v_i\right\rangle=0
             \quad(g\in G_\epsilon,\ v'\in V).
 \tag{BQS.29}
\]
Let $C$ be the closed span of all $\sigma(g)V$.  It contains every
word on $V$: real derived actions are limits of difference quotients
in this closed invariant subspace, and complex matrix units are
complex linear combinations of real directions.
The vector in (BQS.29) lies in $C$ and is orthogonal to every
$\sigma(g^{-1})V$, hence to $C$.  It is zero.
This holds for every coordinate $c$, proving $\mathcal T x=0$.
The argument applies to each $x\in W$ without a finite-dimensional
assumption on $W$ itself.
\end{proof}
% END BQS9

% BEGIN BQS10
\paragraph{BQS10. The universal cone-and-line transfer.}
\label{claim:BQS10}
Here is the full abstract statement used by
\texttt{lowestKType\_tensor\_lift}.  Let $\mathcal S$ be any complex
vector space with compact matrix-unit operators $\mathcal A$ and an
operator $\mathcal B$ commuting with them.  Let
$T:\mathcal S\to M^n$ and $\Theta:\mathcal S\to M_0^n$ be linear,
where $M$ is any complex vector space with an arbitrary operator
family $E$, and $M_0$ is a cyclic matrix-unit module satisfying (BQS.12), with strictly
decreasing $\lambda$.  Let $P\subset M$ be a subspace.  Assume
\[
\begin{aligned}
 T\mathcal B&=\Omega^+_E T,&
 \Theta\mathcal B&=\Omega^+_{\mathcal E}\Theta,\\
 \Theta\mathcal A_{aa}
  &=A^+_{\mathcal E,aa}\Theta+
           \bigl(\Lambda(\lambda)_a-\lambda_a-\rho_a\bigr)\Theta.
\end{aligned}
\tag{BQS.30}
\]
Assume also that every compact highest vector in $P^n$ has a
same-weight highest lift through $T$, and that for every compact-stable
$W\subset\mathcal S$, $\Theta W=0$ implies $TW=0$.
Then every nonzero compact highest vector $f$ of weight $\nu$ in
$P^n\cap\ker(\Omega^+_E-\theta_s^+)$ satisfies
\[
                      \nu-\Lambda(\lambda+e_s)\in Q_+,
 \tag{BQS.31}
\]
and its highest-vector space of weight $\Lambda(\lambda+e_s)$
has dimension at most one.

\begin{proof}
Put $q_s(z)=\prod_{t\ne s}(z-\theta_t^+)$ and
$a_s=q_s(\theta_s^+)\ne0$.
If $Tu=f$ is a highest lift, $u'=q_s(\mathcal B)u$ is still highest,
because $\mathcal B$ commutes with the compact action, and
$Tu'=a_sf$.
BQS4 and the second intertwining in (BQS.30) show that
$\Theta u'$ lies in the selected Euler eigenspace in $M_0^n$.
The same is true of every compact lowering word on $u'$.
If that word lowers the weight by $\delta$, a sum of positive compact
roots, the diagonal relation in (BQS.30) gives its $\Theta$-image
weight
\[
                  \nu-\delta+\lambda+\rho-\Lambda(\lambda).
 \tag{BQS.32}
\]
This identity follows by induction on the word using the compact
matrix-unit commutator; it does not assert that $\Theta$ intertwines
the off-diagonal compact operators.

If $\nu-\Lambda(\lambda+e_s)\notin Q_+$, BQS5 forces every such
image to vanish.  Otherwise a nonzero image would imply
$\nu-\Lambda(\lambda+e_s)-\delta\in Q_+$, and adding $\delta$
would contradict the assumption.  Compact PBW expresses the compact
module generated by the highest vector $u'$ as the span of these
lowering words: commute upper units to the right, where they kill
$u'$, and absorb diagonal units into weights.
This span is compact-stable.  The assumed trace-kernel implication
therefore gives $Tu'=0$, contrary to $a_sf\ne0$.
This proves (BQS.31).

Now take two highest vectors of weight $\Lambda(\lambda+e_s)$,
lift them to $u_1,u_2$, and apply $q_s(\mathcal B)$.
Their two $\Theta$-images have weight $\lambda+\rho+e_s$ in the
selected cyclic-model eigenspace.  BQS5 gives scalars $c_1,c_2$,
not both zero, whose combination has zero image.
Every nonempty compact lowering word on the corresponding highest
combination has model weight $\lambda+\rho+e_s-\delta$ with
$\delta\ne0$ a positive-root sum.  BQS5 and positivity of height
force these images to vanish too.
Apply the trace-kernel implication to its compact lowering span.
The resulting equality
$a_s(c_1f_1+c_2f_2)=0$ gives the required linear dependence.
No finite-dimensionality of $\mathcal S$ or ambient bracket assumptions
beyond the stated compact relations have entered this argument.
\end{proof}
% END BQS10

% BEGIN BQS11
\paragraph{BQS11. The full original standard eigenspace cone and line.}
\label{claim:BQS11}
For every original unitary datum and every $s$, every nonzero compact
highest vector $f\in X_s^+$ of weight $\nu$ has
$\nu-\Lambda(\lambda+e_s)\in Q_+$.
Any two compact highest vectors of weight $\Lambda(\lambda+e_s)$
in $X_s^+$ are linearly dependent.  These are both clauses of the
original \texttt{EigenspaceWeightConeStandard}.

\begin{proof}
Use BSP2 to identify this datum's lowest type isometrically with $T$
and obtain (BQS.24), retaining the given $\lambda,h,\sigma,D,E,V$.
The vector $v_0=\iota^{-1}w_T$ is a compact highest vector of weight
$\Lambda(\lambda)$: differentiate compact intertwining, first in real
directions and then complexify.
Thus $\ell(E_{aa}v)=\Lambda(\lambda)_a\ell(v)$, using the
compact adjoint identity and reality of the weight.

Use BQS6 for the cyclic model, BQS8 for the formal word space, and
BQS9 for its trace-kernel implication.
The two Euler intertwinings in (BQS.30) follow by prepending a letter.
On $[\mathbf Z,v,c]$, the diagonal commutator in the first term of
$\mathcal A_{aa}$ gives
$\mathcal E_{aa}\mathcal E_{\mathbf Z}\Phi_P
 -(\lambda_a+\rho_a)\mathcal E_{\mathbf Z}\Phi_P$.
The $V$ term gives $\Lambda(\lambda)_a\ell(v)$, and the coordinate
term gives the standard tensor summand.
Their sum is exactly the last equality of (BQS.30).
Highest lifts are BQS8.  Therefore BQS10 applies with $P=S$ and
gives both original conclusions.
\end{proof}
% END BQS11

% BEGIN BQS12
\paragraph{BQS12. Actual reversed duality and the whole dual cone statement.}
\label{claim:BQS12}
Both conclusions of BQS11 hold for $X_s^-$ with
$\Lambda(\lambda-e_s)$ in place of $\Lambda(\lambda+e_s)$.
This proves the full original \texttt{EigenspaceWeightConeDual}.
The transport below also preserves arbitrary infinitesimally unitary
target modules, without completing them.

\begin{proof}
Let $\bar a=n-1-a$ and let $j_a\in\{1,-1\}$ be the $a$th diagonal
entry of $J_\epsilon$.  On the same space $D$ set
\[
 \epsilon'_a=\epsilon_{\bar a},\quad
 \lambda'_a=-\lambda_{\bar a},\quad
 E'_{ab}=-j_{\bar a}j_{\bar b}E_{\bar b,\bar a},\quad
 (Uf)_a=j_{\bar a}f_{\bar a}.
 \tag{BQS.33}
\]
This complex-linear signed reversal $U$ is invertible.
The actual new group representation is
$\sigma'(g)=\sigma(R\overline gR)$, where $R e_a=e_{\bar a}$.
The map $g\mapsto R\overline gR$ preserves products, sends
$G_{\epsilon'}$ and $K_{\epsilon'}$ to their original counterparts,
and preserves real one-parameter subgroups.
For a real Lie matrix $X$ the relation
$X^*J_{\epsilon'}+J_{\epsilon'}X=0$ gives
$\overline{X_{ba}}=-j_{\bar a}j_{\bar b}X_{ab}$.
It follows that the complexified derived operators of $\sigma'$ are
exactly $E'$ in (BQS.33).

The matrix-unit relations, Casimir and compact adjoint identities are
preserved.  For example
$(E'_{ab})^*=j_{\bar a}j_{\bar b}E'_{ba}$, the required adjoint sign
for $\epsilon'$, and summing $E'_{ab}E'_{ba}$ cancels both signs and
reindexes the old Casimir.
A highest vector of weight $\nu$ becomes highest of weight
$-\nu\circ\bar{\ }$: if $a<b$, then $\bar b<\bar a$, so
$E'_{ab}$ is a scalar multiple of an original upper unit.
Compact irreducibility and compact stability are unchanged.
The half-root sums satisfy
$\rho_c(\epsilon')=-\rho_c(\epsilon)\circ\bar{\ }$ and likewise for
$\rho_n$, by reversing each root.
Thus
$\Lambda_{\epsilon'}(\lambda')=-\Lambda_\epsilon(\lambda)\circ\bar{\ }$,
$c(\lambda')=c(\lambda)$, and the same original gap and $h_{\rm need}$
hold.  Integrality is preserved because
$-\lambda_{\bar a}-(n-1)/2$ differs from
$-(\lambda_{\bar a}-(n-1)/2)$ by the integer $-(n-1)$.

Direct substitution in (BQS.4) gives
\[
 \Omega^+_{E'}U=U\Omega^-_E,\qquad
 A^+_{E',ab}U
     =U\bigl(-j_{\bar a}j_{\bar b}A^-_{E,\bar b,\bar a}\bigr).
 \tag{BQS.34}
\]
For example in the Euler identity the two coefficients
$j_{\bar c}$ from $U$ cancel $j_{\bar c}^2=1$, leaving
$-j_{\bar d}\sum_cE_{\bar d,c}f_c$.
Generated modules also agree, since (BQS.33) is an invertible
relabeling and rescaling of all matrix units.
Therefore $U$ sends $X_s^-$ to the standard eigenspace for
$(\epsilon',\lambda',\bar s)$, and preserves highest-vector relations
with the just described reversed weights.

Apply BQS11 there.  Its shifted lowest weight is
$-\Lambda_\epsilon(\lambda-e_s)\circ\bar{\ }$.
The map $\beta\mapsto-\beta\circ\bar{\ }$ preserves $Q_+$, because
it sends $\alpha_{ab}$ to $\alpha_{\bar b,\bar a}$.
This proves the cone conclusion; invertibility and linearity of $U$
give the extremal-line conclusion.
The same formula (BQS.33) for operators on any inner-product target
preserves its matrix-unit and adjoint identities by the same checks.
No use of target completeness occurs.
\end{proof}
% END BQS12

\subsection{The lowest type in every nonzero unitary image}

% BEGIN BQS13
\paragraph{BQS13. One highest vector controls every target intertwiner.}
\label{claim:BQS13}
More generally, let $Y$ be any complex inner-product space with a
family of compact operators $A_{ab}$ satisfying the compact adjoint
identities.  Let $X\subset Y$ have finite compact-stable hulls inside
$X$, and assume its highest-vector space of a prescribed real weight
$\nu$ has dimension at most one.
There exists one $y_s\in X$, chosen before any target, such that for
every complex vector space $N$, every family of operators $B_{ab}$
on $N$ and every linear map $\varphi:Y\to N$ intertwining the compact
operators on $X$, every highest vector $w$ of weight $\nu$ in
$\varphi(X)$ is a scalar multiple of $\varphi(y_s)$.
Neither matrix-unit relations on $A$ nor a norm, completeness,
bracket or irreducibility hypothesis on the target are needed.
In particular this holds for either sign with $Y=D^n$,
$X=X_s^\pm$ and $\nu=\Lambda(\lambda\pm e_s)$ for every original
unitary datum.  It supplies the full abstract \texttt{line} and
its original \texttt{line\_V} application.

\begin{proof}
If $w=\varphi(x)$, enclose $x$ in a finite compact-stable
subspace $F\subset X$ as assumed.
The kernel of $\varphi|_F$ is compact-stable, and its orthogonal
complement is compact-stable.  Inverting $\varphi$ on that complement
lifts $w$ to a highest vector of the same weight in $X$.
This proof uses only the intertwining equations on $F$.

If the source highest-vector space is zero, take $y_s=0$; the lifting
argument says that every indicated $w$ is zero.
Otherwise choose any nonzero $y_s$ in it.  The assumed dimension
bound says that this source highest-vector space is one-dimensional.  Every lift is
therefore a scalar multiple of $y_s$, proving the assertion uniformly
for all targets and all maps.
For the stated applications, both tensor actions obey the compact
adjoint identities for the sum inner product on $D^n$.
BQS1 supplies the finite hulls and BQS11 or BQS12 supplies the
one-dimensional source highest space.  Thus the conclusion has
the required original quantifier order, with the same vector for
every target and map.
\end{proof}
% END BQS13

% BEGIN BQS14
\paragraph{BQS14. The full Schmid--Casimir comparison without target completeness.}
\label{claim:BQS14}
Let $N$ be any complex inner-product space and $B$ a matrix-unit
representation with
$B_{ab}^*=j_aj_bB_{ba}$ in the sense of the inner-product identities.
Let $W\subset N$ be a compact irreducible type of highest weight $\nu$.
Put
$\eta=\nu+\rho_c-\rho_n$, and suppose $\eta$ is regular integral,
its noncompact gaps are at least $b$, $p,q\ge1$, and
$h_{\rm need}(\epsilon)<b+2$.
On $\mathfrak p_{\mathbb C}\otimes W$ let
\[
 A_W=\sum_{\epsilon_a=\epsilon_b}
           \operatorname{ad}(E_{ab})\otimes B_{ba}|_W,\qquad
 m(X\otimes w)=B_Xw.
 \tag{BQS.35}
\]
Let $W_-$ be the sum of the negative eigenspaces of this self-adjoint
operator, exactly the original \texttt{Wm}.
If $m(W_-)=0$ and the Casimir has scalar value $c\in\mathbb C$
on $W$, then
\[
                               c=c(\eta).
 \tag{BQS.36}
\]
This is the whole original \texttt{SchmidCasimirStatement}.

\begin{proof}
The first two bounds in (BQS.3) give
\[
 b>npq+1,\qquad b>3n/2,\qquad
 b>16nq_0+8\sqrt{2q_0}.
 \tag{BQS.37}
\]
DC1--DC11 therefore construct the actual unitary model of parameter
$\eta$.  Let $W_0$ be its lowest compact type.  It has highest weight
$\nu$, Casimir $c(\eta)$, and the same negative first-layer relation
by DC8.  Choose an isometric compact intertwiner $i:W\to W_0$ and
put $\Delta=\{(w,iw):w\in W\}$ inside the inner-product direct sum
of $N$ and the model's smooth module.
This direct sum need not be complete.

The compact coupling operator (BQS.35) is natural for compact
intertwiners: $1\otimes i$ intertwines every summand.  Hence the
negative spectral subspace for $\Delta$ is the diagonal copy of the
negative spectral subspace for $W$, also identified with that for
$W_0$.  Multiplication kills it in each coordinate, and therefore
kills it in the direct sum.

Apply the purely algebraic finite-layer argument of DC9 to $\Delta$.
For precision, it requires only this negative first-layer relation
and $b>npq+1$: moving a negative tensor slot to the last position
changes a word by smaller-length terms; a minimal-norm lift of a new
compact type in degree $m$ consequently has every slice in the positive
coupling space.  The compact Casimir identity then gives
\[
 c_K(\tau)\ge c_K(\nu)+2m(b-n+1)\quad(m\ge1).
 \tag{BQS.38}
\]
All lifts and complements in this argument are finite-dimensional.
Since $p,q\ge1$ implies $pq\ge n-1$, the factor $b-n+1$ is positive.
Thus the original compact type $\nu$ occurs in the generated module
only in degree zero and has its one highest line.

Take a nonzero highest vector $(w,iw)$ of $\Delta$.
The total Casimir commutes with the compact action, so its image is
another highest vector of weight $\nu$ in the generated module.
By the preceding uniqueness it is a scalar multiple of $(w,iw)$.
But its two coordinates are $cw$ and $c(\eta)iw$.
As both coordinates are nonzero, the two scalars agree.
This proves (BQS.36) using the same finite-layer comparison as the
source's diagonal-model argument; it assumes neither a completion nor
an actual noncompact group action on $N$.
\end{proof}
% END BQS14

% BEGIN BQS15
\paragraph{BQS15. The exact two-unit gap margin for a candidate compact type.}
\label{claim:BQS15}
Let $\lambda,h$ satisfy the original datum hypotheses, fix $s,t$, and
suppose $\nu$ is compact dominant with
\[
 \nu-\Lambda(\lambda+e_s)\in Q_+,\qquad
 P_j(\nu)\le P_j(\Lambda(\lambda)+e_t)\quad(0\le j\le n).
 \tag{BQS.39}
\]
Then $\eta=\nu+\rho_c-\rho_n$ is regular integral, has noncompact
gaps at least $h-2$, and
\[
 h-2>npq+1,\qquad h_{\rm need}< (h-2)+2,\qquad
 \eta=(\lambda+e_s)+
                 \bigl(\nu-\Lambda(\lambda+e_s)\bigr).
 \tag{BQS.40}
\]

\begin{proof}
Put $\gamma=\nu-\Lambda(\lambda)$ and
$\beta=\nu-\Lambda(\lambda+e_s)\in Q_+$.
Then $\gamma=e_s+\beta$.  The root-cone lower bounds and (BQS.39)
give
\[
        0\le P_j(\gamma)\le1,\qquad
        -1\le\gamma_i=P_{i+1}(\gamma)-P_i(\gamma)\le1.
 \tag{BQS.41}
\]
In particular $|\gamma_a-\gamma_b|\le2$, and every $\gamma_i$ is an
integer.  As $\eta=\lambda+\gamma$, integrality and the noncompact
gap at least $h-2$ follow.

To prove strict decrease it suffices to check adjacent indices $a,a+1$.
For opposite signs,
$\eta_a-\eta_{a+1}\ge h-2>0$.
For equal signs, the half-root counts give
$(\rho_c)_a-(\rho_c)_{a+1}=1$ and
$(\rho_n)_a-(\rho_n)_{a+1}=0$.
Compact dominance gives
$\eta_a-\eta_{a+1}=\nu_a-\nu_{a+1}+1\ge1$.
Finally $h_{\rm need}\ge npq+3$, by (BQS.3), so $h-2>npq+1$.
The other identities in (BQS.40) are the original strict inequality
$h>h_{\rm need}$ and the definition of $\Lambda$.
The margin has not been enlarged.
\end{proof}
% END BQS15

% BEGIN BQS16
\paragraph{BQS16. Every nonzero standard tensor image contains the shifted lowest type.}
\label{claim:BQS16}
Let an original unitary datum and an index $s$ be given, and assume only
that $\lambda+e_s$ is strictly decreasing.
Let $N$ be any complex inner-product space, possibly incomplete, carrying
matrix-unit operators $B$ with the infinitesimal unitary adjoint signs
for $\epsilon$.
For every linear map $\varphi:D^n\to N$ satisfying
\[
       \varphi(A^+_{ab}f)=B_{ab}\varphi(f)
                  \quad(a,b,\ f\in S^n)
 \tag{BQS.42}
\]
and every $x\in X_s^+$ with $\varphi(x)\ne0$, there exists a subspace
$W\subset\varphi(X_s^+)$ which is compact irreducible of highest weight
$\Lambda(\lambda+e_s)$ and on which the Casimir is $c(\lambda+e_s)$.
This is the full original \texttt{QuotientHighestWeightTensorStandard}.
It assumes neither an actual group representation on $N$ nor a chosen
coefficient identity for the source datum.

\begin{proof}
Let $Y=\varphi(X_s^+)$ and let $\Pi=\Pi_s^+$ from BQS3.
BQS1 and (BQS.42) show that $Y$ is Lie-stable and compact-finite.
If $\varphi\Pi$ vanished on $V^n$, its intertwining and the tensor
generation in BQS1 would force it to vanish on all of $S^n$.
This contradicts $\Pi x=x$ and $\varphi(x)\ne0$.
Thus the finite compact-stable image $\varphi\Pi(V^n)$ is nonzero.

Finite-dimensional compact-unitary modules are orthogonal sums of
compact irreducibles.  The elementary argument suffices here:
the commuting self-adjoint diagonal operators have simultaneous real
weight spaces; choose a maximal weight for the compact positive order
to obtain a highest vector; its generated irreducible summand can be
split off orthogonally and the procedure repeated.
Choose such a highest vector in $\varphi\Pi(V^n)$, of weight $\nu_0$.
Lift it through $\varphi\Pi|_{V^n}$ using the finite-dimensional
kernel-complement argument of BQS13.
If coordinate $t$ of its nonzero highest lift is nonzero, that coordinate
is a weight vector in $V$ of weight $\nu_0-e_t$.
Every weight of the compact highest-weight module $V$ differs from
$\Lambda(\lambda)$ by the negative of a sum of positive compact roots:
compact PBW moves all upper units to its highest vector.
Consequently
\[
              P_j(\nu_0)\le P_j(\Lambda(\lambda)+e_t)
                         \quad\hbox{for every }j .
 \tag{BQS.43}
\]

Every nonzero highest vector in $Y$ lifts to one in $X_s^+$ by the
same finite-hull argument.  BQS11 therefore gives
$\nu-\Lambda(\lambda+e_s)\in Q_+$ for every compact highest weight
occurring in $Y$.
Among those satisfying the upper bound (BQS.43), choose one $\nu$
minimizing the nonnegative integer
$\operatorname{ht}(\nu-\Lambda(\lambda+e_s))$.
This collection is nonempty by $\nu_0$.
Let $W\subset Y$ be a compact irreducible subspace of highest weight
$\nu$, contained in a finite compact hull.
Its weight is compact dominant; for a compact simple root this follows
also from
$\|B_{ba}w\|^2=(\nu_a-\nu_b)\|w\|^2\ge0$ on a highest vector.

BQS15 applies, so
$\eta=\nu+\rho_c-\rho_n$ is regular integral and has the exact gap
$h-2$ needed in BQS14.
We claim that the negative first noncompact layer of $W$ is killed
by multiplication.
The finite tensor calculation of DC8 applies to $W$ with parameter
$\eta$.  Its negative coupling space is a sum of compact types
\[
                   \nu-\alpha_{ab},
             \qquad a<b,\quad\epsilon_a\ne\epsilon_b.
 \tag{BQS.44}
\]
To recall why the sign is exact, the compact standard/dual tensor rule
in the two sign blocks gives types $\nu+e_r-e_t$.
Their coupling eigenvalues are
$\eta_r+(\rho_n)_r-\eta_t-(\rho_n)_t+1-n/2$.
The gap $h-2>npq+1$ separates their signs: a negative eigenvalue
requires $r>t$, giving (BQS.44).

If multiplication had a nonzero image on this negative space, its
finite-dimensional compact-unitary image would contain a highest
vector of one of the weights (BQS.44).
It belongs to $Y$, since $Y$ is Lie-stable.
Its partial sums are no larger than those of $\nu$, hence still
satisfy (BQS.43).  Its root-cone difference from
$\Lambda(\lambda+e_s)$ remains in $Q_+$ by the preceding highest-lift
argument, but its height has decreased by $b-a>0$.
This contradicts minimality.  Thus the required Schmid relation holds.

The original Casimir is scalar on all of $S$ by centrality and generation.
BST10's exact tensor identity gives on $X_s^+$
\[
 C_{A^+}=c(\lambda)+n+2\theta_s^+
         =c(\lambda)+2\lambda_s+1=c(\lambda+e_s).
 \tag{BQS.45}
\]
Intertwining transfers this scalar to $Y$, and in particular to $W$.
BQS14 now identifies it with $c(\eta)$.
Writing $\beta=\nu-\Lambda(\lambda+e_s)\in Q_+$, BQS15 gives
$\eta=\lambda+e_s+\beta$, and hence
\[
 0=\|\lambda+e_s+\beta\|^2-\|\lambda+e_s\|^2
   =2\langle\lambda+e_s,\beta\rangle+\|\beta\|^2.
 \tag{BQS.46}
\]
The first term is nonnegative: it is a sum of nonnegative integer
multiples of $(\lambda+e_s)_a-(\lambda+e_s)_b$, $a<b$.
The second is nonnegative and vanishes only when $\beta=0$.
Thus $\nu=\Lambda(\lambda+e_s)$, proving both required properties
of $W$.  Every use of orthogonal decomposition was in a finite-dimensional
subspace of $N$, so target completeness was not assumed.
\end{proof}
% END BQS16

% BEGIN BQS17
\paragraph{BQS17. The full original dual unitary-image theorem.}
\label{claim:BQS17}
Under the same original datum, suppose $\lambda-e_s$ is strictly
decreasing.  For every infinitesimally unitary matrix-unit module
$(N,B)$, every $\varphi:D^n\to N$ intertwining $A^-$ on all of $S^n$,
and every $x\in X_s^-$ with $\varphi(x)\ne0$, there exists
$W\subset\varphi(X_s^-)$ compact irreducible of highest weight
$\Lambda(\lambda-e_s)$ and Casimir $c(\lambda-e_s)$.
This is the entire original \texttt{QuotientHighestWeightTensorDual},
at the same possibly incomplete target scope.

\begin{proof}
Apply the actual reversal BQS12 to the source and set
$B'_{ab}=-j_{\bar a}j_{\bar b}B_{\bar b,\bar a}$ on the same target
$N$.  Its matrix-unit and infinitesimal unitary identities were checked
there.  The map
$\varphi'=\varphi U^{-1}$ intertwines $A^+_{E'}$ with $B'$ on
the whole transformed generated module, by (BQS.34).
It is nonzero on $Ux$.
The relevant standard shift is
\[
            \lambda'+e_{\bar s}
                     =-(\lambda-e_s)\circ\bar{\ },
 \tag{BQS.47}
\]
which is strictly decreasing precisely when $\lambda-e_s$ is.
BQS16 gives the required compact irreducible space $W$ in the same
image $\varphi(X_s^-)$, but expressed using $B'$.
Reversing again returns its highest weight to
$\Lambda_\epsilon(\lambda-e_s)$.
The Casimir is unchanged by the signed transpose reversal and
$c(-(\lambda-e_s)\circ\bar{\ })=c(\lambda-e_s)$.
This proves the original dual conclusion with all source and target
quantifiers preserved.
\end{proof}
% END BQS17

% BEGIN BQS18
\paragraph{BQS18. The original closure and compact-finite assertions.}
\label{claim:BQS18}
Let $\sigma$ be an actual unitary representation of $G_\epsilon$ on
an inner-product space $H$, with derived invariant domain $(D,E)$.
If $W\subset D$ is finite dimensional, stable under the compact
matrix units, and has a real scalar Casimir, then
\[
 C=\overline{U(\mathfrak g_{\mathbb C})W}^{\,H}
                 \quad\hbox{is invariant under every }\sigma(g),
                 \quad g\in G_\epsilon .
 \tag{BQS.48}
\]
Completeness of $H$ is unnecessary for this assertion.
If $H$ is complete and, additionally, $W$ is compact irreducible
of weight $\Lambda(\eta)$, with Casimir $c(\eta)$, and
\[
 \eta\text{ is strictly decreasing},\qquad
 \exists b_0>npq+1:\quad
 |\eta_a-\eta_b|\ge b_0\quad(\epsilon_a\ne\epsilon_b),
 \tag{BQS.49}
\]
then every
finite-dimensional actual $K_\epsilon$-stable subspace of $C$
lies in $U(\mathfrak g_{\mathbb C})W$.
This second assertion is the full original \texttt{data\_kfin\_le}
under \texttt{StdGap}; it assumes no integrality or coefficient identity.

\begin{proof}
Write $F_j$ for the finite span of words of at most $j$ noncompact
matrix units on $W$.  Commuting a compact unit to the right shows
that every $F_j$ is compact stable and that their union is the
generated module.  A finite compact-stable subspace is stable under
the actual compact group.  Indeed, for a compact real Lie algebra
element $X$, the finite-dimensional curve $e^{tE(X)}w$ and
$\sigma(e^{tX})w$ solve the same differential equation with the
same initial value; differentiating
$\sigma(e^{-tX})e^{tE(X)}w$ makes it constant.
Every unitary matrix in each sign block has a skew-Hermitian logarithm,
so this proves stability for all of $K_\epsilon$.
The same argument applies to $F_j$.

Let $c$ be the scalar Casimir.  Centrality extends it from $W$
to every $F_j$.  The finite compact tensor decomposition of
$\mathfrak p_{\mathbb C}^{\otimes j}\otimes W$ bounds its compact
highest weights by $A+\sqrt2j$ for a fixed $A$.
The energy identity
$L=2C_K-C_G$, obtained by summing the squares of a real orthonormal
compact and noncompact basis with the infinitesimal unitary adjoint
signs, therefore bounds the norm of $E(X)$ on $F_j$ by
$C_X(A'+j)$ for fixed $A'\ge1$ and $C_X$.
Iterating, for $v\in F_k$,
\[
 \|E(X)^m v\|
 \le C_X^m\prod_{r=0}^{m-1}(A'+k+r)\,\|v\|.
 \tag{BQS.50}
\]
For $|t|C_X<1$ the Taylor series converges.
The integral Taylor remainder for
$t\mapsto\sigma(e^{tX})v$ has norm at most
$|t|^m\|E(X)^m v\|/m!$; unitarity gives this bound at every
base point.  Its successive majorants have ratio tending to
$|t|C_X<1$.  The remainder tends to zero, so the series has
the actual orbit as its sum, also if $H$ is incomplete.
All its partial sums lie in the generated module, and thus
$\sigma(e^{tX})v\in C$ for these $t$.
Continuity of the bounded operator $\sigma(e^{tX})$ extends
this inclusion from the dense generated module to $C$.
Subdividing $t=1$ proves stability under $e^X$.
Polar decomposition $g=k e^X$, $X\in\mathfrak p$, together with
the compact stability already proved gives (BQS.48).
Applying it to $g^{-1}$ gives equality.

For the second assertion write $b_0$ for the real bound in
(BQS.49).  If one sign block is empty, all matrix units are compact,
the generated module is $W$, and $C=W$ is already finite dimensional.
Otherwise $b_0>npq+1$ implies $b_0-n+1>0$.
The finite Casimir argument in DC8 supplies the negative first-layer
Schmid relation at exactly this bound.  DC9 then gives, on the
orthogonal new part of $F_j$, $j\ge1$,
\[
 c_K(\nu)\ \ge\ c_K(\Lambda(\eta))
                         +2j(b_0-n+1)
 \tag{BQS.51}
\]
for each nonzero compact highest vector of weight $\nu$.
Neither argument uses integrality of $\eta$: the compact module
$W$ is already supplied and its finite-dimensional unitary
decomposition is the only compact input.
Consequently every vector of a fixed compact type $\nu$ in the
generated module lies in some one finite layer $F_{J(\nu)}$.
To see this for arbitrary word length, decompose a vector
orthogonally through the nested finite compact-stable layers.
Each nonzero new component has the same type and its layer index
is bounded by (BQS.51).

The actual compact character projector onto $\nu$ is bounded on
$H$: it is the integral over probability Haar of the continuous
unitary operators multiplied by the finite-dimensional character
and its dimension.  On every finite compact hull, finite
Schur orthogonality identifies its range with the $\nu$-isotypic
subspace.  Its image on the dense generated module is therefore
contained in $F_{J(\nu)}$, a closed finite-dimensional subspace.
Continuity gives the same inclusion on $C$.
Any finite-dimensional actual compact-stable subspace of $C$ is
the direct sum of finitely many such compact types, so it lies
in the generated module.  This proves the stated full weak-gap
assertion without inserting the stronger constructive gap.
\end{proof}
% END BQS18

% BEGIN BQS19
\paragraph{BQS19. The full universal-image closure lemma.}
\label{claim:BQS19}
Let $Y$ be any complex vector space, $A_{ab}\in\operatorname{End}Y$
any family of linear operators, and $G,X\subset Y$ any subspaces.
No matrix-unit relations on $A$, stability of $G$ or $X$, or inclusion
$X\subset G$ is assumed here.
Fix a real parameter $\eta$ and write
$\mathcal D=\operatorname{smD}(r_{\rm reg})$ for the smooth domain
in the actual right regular $L^2(G_\epsilon)$ representation, with
derived operators $B$.
Call $\psi:Y\to\mathcal D$ admissible when
\[
               \psi(A_{ab}y)=B_{ab}\psi(y)
                            \quad(y\in G,\ a,b).
 \tag{BQS.52}
\]
Assume the following two universal properties.
For every admissible $\psi$, every $x\in X$ with $\psi(x)\ne0$
gives a compact irreducible $W_\psi\subset\psi(X)$ of weight
$\Lambda(\eta)$ and Casimir $c(\eta)$.
There is one specified $y_s\in Y$ such that, for every admissible
$\psi$, every highest vector of weight $\Lambda(\eta)$ in $\psi(X)$
is a scalar multiple of $\psi(y_s)$.
Given an admissible $\varphi$ and a compact irreducible
$W\subset\mathcal D$ of this weight and Casimir with
$\varphi(y_s)\in W$, one has
\[
       \varphi(X)\subset
       \overline{U(\mathfrak g_{\mathbb C})W}^{\,L^2}.
 \tag{BQS.53}
\]
This is the original \texttt{image\_le\_closGen}, without a gap
assumption and without an assumption that $W\subset\varphi(X)$.

\begin{proof}
Let $C$ be the closure on the right.  BQS18 makes it invariant
under the actual right regular representation, and unitarity and
invariance under inverses make $C^\perp$ invariant as well.
Its bounded orthogonal projection $P$ commutes with that representation.
If $v$ is smooth, its orbit is locally the restriction of a smooth
ambient map; composing that map with $P$ proves that $Pv$ is smooth.
Differentiating $P r(g)v=r(g)Pv$ along each real one-parameter
subgroup and complexifying gives $PB_{ab}v=B_{ab}Pv$.
Thus $\psi=P\varphi$, viewed as a map to $\mathcal D$, is admissible.
As $\varphi(y_s)\in W\subset C$, one has $\psi(y_s)=0$.

If $\psi(x)\ne0$ for some $x\in X$, the universal image property
gives a nonzero compact irreducible $W_\psi\subset\psi(X)$.
Take its nonzero highest vector $w$.
The universal line property says $w=a\psi(y_s)=0$, a contradiction.
Therefore $P\varphi(X)=0$.  Since $C$ is closed, $\ker P=C$,
which is exactly (BQS.53).
Completeness has been used only for the actual $L^2$ orthogonal
projection; the source $Y$ remains arbitrary.
\end{proof}
% END BQS19

% BEGIN BQS20
\paragraph{BQS20. The entire original generic quotient theorem.}
\label{claim:BQS20}
Use the arbitrary source $Y,A,G,X$, parameter $\eta$, target
$\mathcal D$ and admissibility condition (BQS.52) of BQS19.
Assume now $X\subset G$ and \texttt{StdGap} (BQS.49) for $\eta$.
Retain the universal nonzero-image property of BQS19, and assume
there exists \emph{one} $y_s\in X$ with its universal highest-line
property for every admissible map.
Suppose every $x\in X$ lies in a finite-dimensional compact
matrix-unit-stable subspace $S_x\subset X$.
Let $Z\subset X$ be stable under every $A_{ab}$, and let $\varphi$
be admissible and satisfy the exact saturation condition
\[
 y\in X,\quad z\in Z,\quad \varphi(y)=\varphi(z)
                       \quad\Longrightarrow\quad y\in Z.
 \tag{BQS.54}
\]
For every $x\in X\setminus Z$ there exist subspaces $Y'\subset X$
and $W\subset\mathcal D$ such that
\[
 \varphi(Y')=W,\qquad Y'\cap Z=0,\qquad
 W\text{ is compact irreducible of weight }\Lambda(\eta),
 \qquad C_B|_W=c(\eta),
 \tag{BQS.55}
\]
and $W$ is stable under every element of the actual $K_\epsilon$.
No relations on $A$, complete norm on $Y$, global intertwining
outside $G$, or full Lie stability of $X$ has been added.

\begin{proof}
If $\varphi(x)=0=\varphi(0)$, (BQS.54) would give $x\in Z$.
Thus the universal image property gives a compact irreducible
$W\subset\varphi(X)$ of the required weight and Casimir.
Its nonzero highest vector has the form $a\varphi(y_s)$ by the
universal line property.  Here $a\ne0$, so $\varphi(y_s)\in W$.
BQS19 gives $\varphi(X)\subset C=\overline{U(\mathfrak g)W}$.

The image $\varphi(S_x)$ is finite dimensional and compact
matrix-unit-stable: $S_x\subset X\subset G$ permits (BQS.52)
on every vector of $S_x$.  The finite-dimensional group
argument in BQS18 makes it stable under the actual compact
group.  The weak-gap conclusion of BQS18 therefore implies
\[
                       \varphi(x)\in U(\mathfrak g)W.
 \tag{BQS.56}
\]
Put $N_0=\varphi(Z)$.  For $z\in Z$, the inclusion
$z\in X\subset G$ and (BQS.52) give
$B_{ab}\varphi(z)=\varphi(A_{ab}z)\in N_0$.
Hence $W\cap N_0$ is compact stable, so compact irreducibility
makes it either zero or $W$.
The second alternative would imply
$U(\mathfrak g)W\subset N_0$.
By (BQS.56), $\varphi(x)=\varphi(z)$ for some $z\in Z$,
contradicting (BQS.54).  Thus $W\cap N_0=0$.

Let $Y_0=X\cap\varphi^{-1}(W)$.
Since $W\subset\varphi(X)$, the restriction
$\varphi:Y_0\to W$ is surjective.
Choose an algebraic complement $Y'$ to its kernel in $Y_0$.
Then $\varphi|_{Y'}$ is an isomorphism onto $W$.
If $y\in Y'\cap Z$, its image lies in $W\cap N_0=0$,
and injectivity gives $y=0$.
This proves all of (BQS.55).  Finally $W$ is finite dimensional
and stable under the compact derived operators in the actual
smooth right regular representation; the logarithm and orbit
argument in BQS18 proves its actual $K_\epsilon$ stability.
The complement is an algebraic choice and imposes no unwanted
continuity requirement on the original source.
\end{proof}
% END BQS20

\subsection{The original smooth maps and the actual period kernel}

For the last three claims let $n=2r+k+1$, $p=k+r$, $q=r+1$.
The original step hypotheses are
\[
 \max(3n/2,1+npq)<h-2,\qquad
 4\max(2nq_0,2rn^2)<h/(8\sqrt{2q_0})-1.
 \tag{BQS.57}
\]
These are exactly \texttt{StepGap}; BSP1 proves $q_0\ge1$
from these inequalities, including the original zero-division convention.
Let $\lambda$ be Harish--Chandra with noncompact gaps at least $h$.
We retain every original unitary datum with this parameter and, in the
last two claims, the supplied $T,\mu,\iota$ with
$\mu=\lambda+\rho$ and the actual coefficient identity (BSP.4).

The Bessel data are arbitrary original data:
$u\in\mathbb R^k$, a middle vector $z$ with $T_z=-z^*H_{pk}z>0$,
$c\in(\mathbb C^\times)^r$, a complete complex Hilbert space $V_\pi$,
a representation $\pi$ satisfying the full
$\texttt{IsHCTypeW}(k,r,z,u,\pi)$, and ambient matrix measures
$\mu_W,\mu_N$ satisfying the original weak
$\texttt{IsHaarOn}(H_z,\mu_W)$ and
$\texttt{IsHaarOn}(N,\mu_N)$.
Here $H_z=\texttt{UWset}(k,r,z)$, $N=\texttt{Nset}(k,r)$,
and $\psi_{z,c}$ is the actual character in (BNV.1).
Let $P$ be any invertible Cayley matrix
$P^*J_\epsilon P=H_{pk}$.  Euclidean unitarity is not required.
Put $x_1=x_J(P,z)=P\,x_{\rm One}(z)$, with the prescribed Bessel
vector $x_{\rm One}$, and fix any
$x_0\in G_\epsilon$, $w\in V$, and $\xi\in V_\pi$.
For a tuple $f\in D^n$, the exact two maps on ambient matrices are
\[
 \begin{split}
 L^+(f)(g)&=\sum_c
  (x_1^*J_\epsilon ge_c)\,
       \langle w,\sigma(x_0^{-1}g)f_c\rangle,\\
 L^-(f)(g)&=\sum_c
  (e_c^t g^{-1}x_1)\,
       \langle w,\sigma(x_0^{-1}g)f_c\rangle .
 \end{split}
 \tag{BQS.58}
\]
Off the group, the inverse has the same total matrix-inverse convention
as the source.  These are its \texttt{tMap} and \texttt{tMapD},
not equivalence classes of functions.
For either scalar function $F$ the actual vector-valued period is
\[
 \mathcal P_F(g)=
 \int F(P h' n' gP^{-1})\,
       \overline{\psi_{z,c}(n')}\,\pi(h'^{-1})\xi\,
             d(\mu_W\otimes\mu_N)(h',n'),\qquad g\in G_H=U(H_{pk}).
 \tag{BQS.59}
\]
The integral is on the two ambient matrix spaces; the weak Haar
predicates give support on $H_z$ and $N$ almost everywhere.
Every choice in this paragraph, including $\xi$, remains fixed
through the quotient construction.

% BEGIN BQS21
\paragraph{BQS21. The smooth lift and the exact simultaneous period kernel.}
\label{claim:BQS21}
First, in arbitrary rank and signature, let $Y$ be any complex vector
space, let $A_{ab}$ be any endomorphisms, let $G\subset Y$ be stable
under them, and let $L:Y\to(\operatorname{Mat}_n(\mathbb C)\to\mathbb C)$
be linear.
Suppose for every $f\in G$ that $Lf$ is smooth on $G_\epsilon$,
all its derivative words are in $L^2$, and
\[
                   E_{ab}(Lf)=L(A_{ab}f)\quad\hbox{on }G_\epsilon.
 \tag{BQS.60}
\]
Then there is a linear map $\Phi:Y\to\mathcal D$ such that on $G$
\[
        \Phi(f)=\operatorname{toL2}(Lf),\qquad
        \Phi(A_{ab}f)=B_{ab}\Phi(f).
 \tag{BQS.61}
\]
This is the entire original \texttt{exists\_lin}.

For the original data (BQS.57)--(BQS.59), apply this to
$Y=D^n$, $G=S^n$, $A=A^\pm$, $L=L^\pm$.
For each $s$, the subspace
\[
 Z_s^\pm=\{y\in X_s^\pm:
                     \mathcal P_{L^\pm(y)}(g)=0
                           \text{ for every }g\in G_H\}
 \tag{BQS.62}
\]
is stable under every $A^\pm_{ab}$ and satisfies the exact
saturation condition
\[
 y\in X_s^\pm,\quad z'\in Z_s^\pm,\quad \Phi(y)=\Phi(z')
                      \quad\Longrightarrow\quad y\in Z_s^\pm .
 \tag{BQS.63}
\]
If the supplied tuple $x\in X_s^\pm$ has a nonzero period
at some $g\in G_H$, then $x\notin Z_s^\pm$.

\begin{proof}
On $G$, every $\operatorname{toL2}(Lf)$ is an actual $L^2$ class,
so the operation is linear there; the source's value zero outside
the $L^2$ domain never enters.  These classes are smooth vectors.
Here are the analytic details of that last assertion.
For a real Lie algebra element $X$, all derivatives $R_X^mLf$
are finite complex linear combinations of the given word derivatives
and belong to $L^2$.
Right translations are strongly continuous on $L^2$:
it holds first for continuous compactly supported functions by
uniform continuity on a compact neighborhood and then for all
$L^2$ functions by density and unitarity.
The pointwise fundamental theorem of calculus for $Lf$ therefore
becomes the Bochner identity
\[
 r(e^{tX})[Lf]-[Lf]
                =\int_0^t r(e^{sX})[R_XLf]\,ds
                   \quad\hbox{in }L^2 .
 \tag{BQS.64}
\]
The integrand is strongly continuous and has constant $L^2$ norm;
on finite intervals Fubini and Cauchy--Schwarz identify its
Bochner integral with the pointwise integral.
Dividing by $t$ gives the derived vector $[R_XLf]$.
Repeating for every mixed derivative, and using coordinates which
are finite products of exponentials of a real Lie basis, makes
the orbit map $C^\infty$.  A local smooth retraction from ambient
matrices to the Lie group supplies the exact local ambient smooth
extensions used by \texttt{smD}, as in BSP5.
Thus the map $f\mapsto[Lf]$ on $G$ lands in $\mathcal D$ and
(BQS.60) gives its intertwining.
Choose an algebraic projection $q:Y\to G$ with $q|_G=1$ and set
$\Phi(f)=[L(qf)]$.  This proves (BQS.61) on all the originally
required vectors; no continuity or relations on $A$ were assumed.

For the actual maps, BQS1 gives invariance of $S^n$ and of
$X_s^\pm$, and BST5 gives (BQS.60), smoothness, and all word
estimates with the unchanged exponent
$1-h/(8\sqrt{2q_0})$.
Thus the generic construction applies.
BSP5 and BNV1 supply absolute Bochner integrability of every
period in (BQS.59), for every such word and every $g\in G_H$.
Consequently period addition and scalar multiplication are genuine
integral identities.  They make (BQS.62) a linear subspace.

The exact differentiation identity of BSP5 identifies the period
of $E_{ab}(L^\pm y)$ with the corresponding complexified right
derivative, in the Cayley-transported direction, of
$g\mapsto\mathcal P_{L^\pm y}(g)$.
If the latter vanishes everywhere on $G_H$, each real tangential
derivative vanishes, and then so does its complexification.
Together with (BQS.60) and invariance of $X_s^\pm$, this proves
$A^\pm_{ab}Z_s^\pm\subset Z_s^\pm$.
There is no integration by parts against a different character or
a different measure in this step.

Finally, $\Phi(y)=\Phi(z')$ says that the two continuous functions
$L^\pm y$ and $L^\pm z'$ have the same $L^2$ class on
$G_\epsilon$.  Haar measure has full support, so a nonzero value
of their continuous difference would give a nonempty open set
of positive measure on which its absolute value is bounded below.
Their difference is therefore zero at every point of $G_\epsilon$.
For almost every $(h',n')$ in the given ambient product measure,
$h'\in H_z$, $n'\in N$, and hence
$P h'n'gP^{-1}\in G_\epsilon$ for every $g\in G_H$.
The two actual integrands in (BQS.59) agree almost everywhere.
Thus their periods agree for every $g$, proving (BQS.63).
The final assertion is precisely the negation of the defining
condition in (BQS.62).
\end{proof}
% END BQS21

% BEGIN BQS22
\paragraph{BQS22. Full original \texttt{QuotPStatement} and \texttt{QuotMStatement}.}
\label{claim:BQS22}
Take every original datum in (BQS.57)--(BQS.59), every index $s$,
and an arbitrary tuple $x\in X_s^\pm$.
Assume $\lambda'=\lambda\pm e_s$ is Harish--Chandra, its
noncompact gaps are at least the \emph{same} $h$, and
$\mathcal P_{L^\pm(x)}(g)\ne0$ for some $g\in G_H$.
There exist a subspace $W_f$ of complex-valued functions on
ambient matrices and a subspace $W'\subset\mathcal D$ with all
of the following original properties:
\[
 \begin{array}{ll}
 {\rm (a)}&
  \text{Every }F\in W_f\text{ equals }L^\pm(y)
                        \text{ for some }y\in X_s^\pm;\\
 {\rm (b)}&
  v\in W'\ \Longleftrightarrow\
       \exists F\in W_f:\ v=\operatorname{toL2}(F)
                         \quad(v\in\mathcal D);\\
 {\rm (c)}&
  W'\text{ is compact irreducible of highest weight }
                               \Lambda(\lambda');\\
 {\rm (d)}& C_B v=c(\lambda')v\quad(v\in W');\\
 {\rm (e)}&
  r_{\rm reg}(k')W'=W'\quad\text{for every actual }k'\in K_\epsilon;\\
 {\rm (f)}&
  F\in W_f,\quad \mathcal P_F(g)=0\text{ for every }g\in G_H
               \quad\Longrightarrow\quad F|_{G_\epsilon}=0.
 \end{array}
 \tag{BQS.65}
\]
The sign $+$ gives the full original \texttt{QuotPStatement};
the sign $-$ gives the full original \texttt{QuotMStatement}.
The tuples are arbitrary elements of the original Euler eigenspaces,
not only the tuples produced by a complementary minor.

\begin{proof}
BSP1 turns (BQS.57) into $p,q\ge1$ and
$h>h_{\rm need}(\epsilon)$, so the supplied unitary datum is in
the exact scope of BQS1--BQS17.
Also $h>npq+1$; the hypotheses on $\lambda'$ therefore give
the weak \texttt{StdGap} (BQS.49) used by the generic quotient
theorem.  The rank-zero impossibility is supplied by BSP1,
not by an extra restriction on the quantified $k,r$.

Use BQS20 with
\[
 Y=D^n,\quad G=S^n,\quad X=X_s^\pm,\quad
 A=A^\pm,\quad Z=Z_s^\pm,\quad \eta=\lambda',
 \tag{BQS.66}
\]
and with the smooth lift $\Phi$ supplied by BQS21.
Here $X\subset G$, and BQS1 supplies every finite compact hull.
BQS21 supplies stability and saturation of $Z$, admissibility
of $\Phi$, and $x\notin Z$.
For \emph{every} admissible $\psi:Y\to\mathcal D$, the target
$\mathcal D$ has its actual infinitesimally unitary matrix-unit
action.  It need not be complete in its inherited norm.
BQS16, respectively BQS17, therefore supplies the universal
nonzero-image property required by BQS20 at exactly this target.
BQS13 supplies one $y_s\in X$ before the choice of $\psi$ and
the required universal highest-line property.
Thus every hypothesis of BQS20 holds with the same $G,X,Z,\Phi$.

Let $Y'\subset X$ and $W$ be its conclusion, and define
\[
                         W_f=L^\pm(Y'),\qquad W'=W.
 \tag{BQS.67}
\]
Property (a) is immediate.
If $v\in W'=\Phi(Y')$, choose $y\in Y'$ with $v=\Phi(y)$;
BQS21 gives $v=\operatorname{toL2}(L^\pm y)$.
Conversely, if $F=L^\pm y$ with $y\in Y'$ and
$v=\operatorname{toL2}(F)$, the same identity gives $v=\Phi(y)\in W'$.
This proves both directions of the exact image statement (b).
Properties (c), (d), and (e) are exactly the three representation
conclusions of BQS20.

For (f), write $F=L^\pm y$ with $y\in Y'$.
If every prescribed period vanishes, (BQS.62) says $y\in Z$.
Since $Y'\cap Z=0$, one has $y=0$ and hence $F=0$ as a function
on all ambient matrices, which in particular gives the required
pointwise vanishing on $G_\epsilon$.
All measures, character, Cayley matrix, $\xi$, $w$, and $x_0$
in this proof were the original supplied ones.
The same proof with the already established dual image and
line theorems proves every field for the minus sign.
\end{proof}
% END BQS22

% BEGIN BQS23
\paragraph{BQS23. Full one-weight propagation and the original endomorphism nonvanishing.}
\label{claim:BQS23}
The complete original \texttt{StepStatement} now holds.
Consequently the complete original \texttt{BesselNVGap} holds:
for every even $k\ge2$, every $r\ge0$, all original
$z,c,u,v,b$ satisfying the Harish--Chandra, interlacing,
half-integrality and \texttt{GapBig} conditions of (BNV.1)--(BNV.3),
every representation $\sigma$ satisfying
$\texttt{IsDSHpk}(k,r,\epsilon^{48},\lambda,\sigma)$,
every $\pi$ satisfying $\texttt{IsHCTypeW}(k,r,z,u,\pi)$,
every original lowest-type unit vector $w$, and every original
weak Haar pair, there exist $x_0,x\in G_H$ such that
\[
 \int
 \langle\sigma(x_0)w,\sigma(h'n')\sigma(x)w\rangle\,
    \overline{\psi_{z,c}(n')}\,\pi(h')^{-1}\,
                  d(\mu_W\otimes\mu_N)(h',n')\ne0
       \quad\text{in }\operatorname{End}_{\mathbb C}(V_\pi).
 \tag{BQS.68}
\]
The original \texttt{ConcreteSettingExists} also follows on the
same supplied \texttt{Forms0} scope as BWP24.

\begin{proof}
BQS22 discharges both exact quotient premises retained in BSE9.
For clarity, the one-weight argument uses the already proved
EvalP and EvalM to turn a nonzero source coefficient period
into a nonzero period of the respective complementary-minor
tuple.  BST's complementary Euler identity places that tuple
in $X_s^+$ or $X_s^-$.
BQS22 produces its actual function quotient with all six fields
(BQS.65), while BST5 supplies the inherited smoothness and
unchanged decay exponent for every function in it.
The full \texttt{FinStatement} proved in BSP, with the full
\texttt{IntSStatement} of BNV1, therefore produces the shifted
original $\texttt{BesselNZ}$.
The exact evaluation factors in BSE7--BSE8 give precisely the
original exclusions $u_j\ne\lambda_s\pm\tfrac12$.
Both signs are included, so this proves the whole
\texttt{StepStatement}, as assembled in BST11 and BSE9.

BNV23 now has its last premise supplied.
Explicitly, BNV7 chooses an admissible real gap before the finite
walk; BNV18 gives a nonzero period at a sufficiently far deformed
tail parameter; BNV6 walks back using the two step directions
just proved.  BNV21 transports the resulting existential
nonvanishing to every original datum, and BNV22 returns it to
the literal endomorphism-valued integral (BQS.68).
That return includes the independently existential Cayley
witnesses, arbitrary lowest unit vector, arbitrary original
nondegenerate character and Haar normalization, and the
actual order $h'n'$ of the two group factors.
Thus (BQS.68) is the prescribed operator integral, not merely
the existence of an abstract branching functional.

Finally BWP24 consumes exactly this original
\texttt{BesselNVGap}, constructs and normalizes the original
wave packets on the same adelic data, and joins the reviewed
ASG and AMD setting and measure producers.
It gives \texttt{ConcreteSettingExists} for each supplied
\texttt{Forms0} datum.  That premise is retained exactly:
this invocation does not construct the ambient form record.
\end{proof}
% END BQS23

\paragraph{Scope of this completion.}
This chapter completes the original Euler polynomials, K2 cone and
extremal line, K3 image types, both actual smooth quotients,
both one-weight propagations, and the original Bessel nonvanishing
theorem.  Its supplied-form setting consequence is conditional only
on the inherited, explicit \texttt{Forms0} input.
The ambient form supplier, orbital/fiber comparison, Siegel--Weil
comparison, local uniqueness and the downstream global spectral
identifications remain separate producers.
In particular this chapter does not close the whole selected
\texttt{GLOBAL\_SPECTRAL} obligation or assert the unrestricted
RH or Dirichlet-GRH endpoints.
